\chapter{DIZAJN I IMPLEMENTACIJA RAZVIJENOG ASEMBLERSKOG PREVODIOCA}
\label{ch:implementacija}

\napomena{Centralno poglavlje i najveći doprinos rada, 15--20 strana.
Svaka sekcija odgovara jednom modulu u \texttt{src/asm/}. Prikazivati kratke isječke koda,
ne cijele fajlove (cijeli kod je u Prilogu~\ref{sec:prilog-kod}).}

\section{Arhitektura softverskog rješenja}
\label{sec:arhitektura-rjesenja}
\napomena{Dijagram toka: \texttt{.asm} → Lexer → Parser → Pass~1 (tabela simbola, adrese)
→ Pass~2 (kodiranje) → \texttt{.bin}/\texttt{.hex}. Tabela modula i odgovornosti.}

\section{Izbor programskog jezika i razvojnih alata}
\label{sec:alati}
\napomena{Python 3.12 (čitljivost, standardna biblioteka, bez spoljnih zavisnosti),
pytest, Git/GitHub, VS Code/Cursor, \LaTeX. Zašto bez generatora parsera.}

\section{Implementacija leksičkog analizatora}
\label{sec:impl-lexer}
\napomena{\texttt{lexer.py}: vrste tokena, prolaz znak po znak, lokacije (linija/kolona),
sažimanje praznih linija u jedan NEWLINE, dekodiranje brojeva i stringova, greške.}

\section{Implementacija parsera}
\label{sec:impl-parser}
\napomena{\texttt{parser.py}: rekurzivno spuštanje po gramatici iz Listinga~\ref{lst:ebnf};
čvorovi \texttt{LabelDef}, \texttt{Directive}, \texttt{Instruction}.}

\section{Tabela simbola}
\label{sec:impl-simboli}
\napomena{\texttt{symbols.py}: struktura, definisanje i pretraga, detekcija duplikata.}

\section{Prevođenje instrukcija u mašinski kod}
\label{sec:impl-kodiranje}
\napomena{\texttt{encode.py}: tabela opkodova, funkcije za formate, provjera opsega
neposrednih vrijednosti, PC-relativni pomjeraji za grane i skokove, pseudo-instrukcije.}

\section{Generisanje izlaznog fajla}
\label{sec:impl-izlaz}
\napomena{\texttt{output.py}: binarni (little-endian) i heksadecimalni zapis, opciono listing.}

\section{Obrada sintaksnih i semantičkih grešaka}
\label{sec:impl-greske}
\napomena{\texttt{errors.py}: hijerarhija grešaka, format \texttt{fajl:linija:kolona: vrsta: poruka},
primjeri poruka.}

\section{Testiranje funkcionalnosti asemblera}
\label{sec:impl-testiranje}
\napomena{Jedinični testovi po modulu, integracioni testovi nad primjerima, testovi grešaka;
broj testova i pokrivenost.}
